\originalchapter{8.323 作业10题解}
题源:MIT OCW,8.323,2023年春季,2023年4月18日发布.\href{https://ocw.mit.edu/courses/8-323-relativistic-quantum-field-theory-i-spring-2023/mit8_323_s23_pset10.pdf}{原作业PDF};\href{https://ocw.mit.edu/courses/8-323-relativistic-quantum-field-theory-i-spring-2023/mit8_323_s23_pset_10sol.pdf}{官方解PDF}.以下为官方解中文重编与独立补充,并非MIT官方中文版;MIT OCW来源,CC BY-NC-SA 4.0.

原资料用$(-+++)$,本书用$(+---)$;$\Gamma^\mu=\ii\gamma^\mu$、$\bar\psi_{\rm o}=\ii\bar\psi$,原$\e^{\ii p\cdot_{\rm o}x}=\e^{-\ii p\cdot x}$.二分量记号为$\sigma^\mu=(1,\boldsymbol\sigma)$、$\bar\sigma^\mu=(1,-\boldsymbol\sigma)$,与作业9相同,$p^0=E_p>0$、$p^2=m^2$.

\originalsection{原题1}
\begin{exercise}
考虑$m=0$的Dirac作用量.(a) 证明$\psi\mapsto\e^{\ii\alpha\gamma^5}\psi$是不变变换;(b) 构造其Noether流;(c) 求质量项$m\bar\psi\psi$的变换,判断是否不变.
\end{exercise}
\begin{solution}
(a) $\gamma^5$Hermitian且与$\gamma^0$反对易,故$\bar\psi\mapsto\bar\psi\e^{\ii\alpha\gamma^5}$.对常数$\alpha$,用$\e^{\ii\alpha\gamma^5}\gamma^\mu=\gamma^\mu\e^{-\ii\alpha\gamma^5}$可得
$\ii\bar\psi\e^{\ii\alpha\gamma^5}\gamma^\mu\e^{\ii\alpha\gamma^5}\partial_\mu\psi=\ii\bar\psi\gamma^\mu\partial_\mu\psi$.

(b) 可用局部参数检验:$\delta\psi=\ii\alpha(x)\gamma^5\psi$使$\delta\Lag=-(\partial_\mu\alpha)\bar\psi\gamma^\mu\gamma^5\psi$,因此按常用整体符号定义轴流
$j_5^\mu=\bar\psi\gamma^\mu\gamma^5\psi$.在自由方程上$\partial_\mu j_5^\mu=0$.Noether流的整体负号可随生成元约定改变.本题是自由经典/自由量子理论,未引入与规范场耦合后的轴反常.

(c) $m\bar\psi\psi\mapsto m\bar\psi\e^{2\ii\alpha\gamma^5}\psi$
$=m\cos(2\alpha)\bar\psi\psi+\ii m\sin(2\alpha)\bar\psi\gamma^5\psi$,不是一般连续旋转下的不变量.质量非零时直接用方程得$\partial_\mu j_5^\mu=2\ii m\bar\psi\gamma^5\psi$.指数是夹在两个旋量之间的矩阵,不能把它误写成乘在已完成收缩的标量外面的矩阵.
\end{solution}

\originalsection{原题2}
\begin{exercise}
对二分量复旋量$\chi=\psi_L$的Majorana理论
\[
 \Lag_L=\ii\chi^\dagger\sigma^\mu\partial_\mu\chi
 -\frac m2(\chi^T\sigma^2\chi+\chi^\dagger\sigma^2\chi^*),
\]
(a) 写等时正则量子化关系;(b) 写经典运动方程,按$\chi(x)=u(p)\e^{-\ii px}+v(p)\e^{\ii px}$求静止系完整解基;(c) 验证任意动量的完整解基
\[
 u_s(p)=\sqrt{p\cdot\bar\sigma}\,\zeta_s,\qquad
 v_s(p)=-\sqrt{p\cdot\bar\sigma}\,\sigma^2\zeta_s,
 \quad\sigma^3\zeta_\pm=\pm\zeta_\pm,
\]
这是原题$\sqrt{-p\cdot_{\rm o}\bar\sigma}$的转换;(d) 写量子场模式展开;(e) 定义真空并构造正确归一化的单粒子态,与Dirac粒子比较.
\end{exercise}
\begin{solution}
(a) 动能一阶时间导数给$\pi_\chi=\ii\chi^\dagger$;消去共轭场的二类约束后的独立正则关系为
\[
 \{\chi_a(t,\mathbf x),\chi_b^\dagger(t,\mathbf y)\}=\delta_{ab}\delta^{(3)}(\mathbf x-\mathbf y),
 \quad\{\chi_a,\chi_b\}=\{\chi_a^\dagger,\chi_b^\dagger\}=0.
\]
对费米场不能把这些写成普通对易子.

(b) 方程及其共轭为$\ii\sigma^\mu\partial_\mu\chi-m\sigma^2\chi^*=0$与$-\ii\partial_\mu\chi^\dagger\sigma^\mu-m\chi^T\sigma^2=0$.对给定正能量壳动量,分离两种相位得
\[
 (p\cdot\sigma)u-m\sigma^2v^*=0,\qquad
 (p\cdot\sigma)v+m\sigma^2u^*=0.
\]
静止时$p\cdot\sigma=m$,所以$v=-\sigma^2u^*$.取实正交基$\zeta_+=(1,0)^T$、$\zeta_-=(0,1)^T$,得到$u_s(\mathbf0)=\sqrt m\zeta_s$、$v_s(\mathbf0)=-\sqrt m\sigma^2\zeta_s$,$u_s^\dagger u_r=m\delta_{sr}$.对一般复线性组合,$v$系数是$u$系数的共轭:Majorana经典解空间是实线性空间,不能把$u,v$视为两组独立复幅度.

(c) 置$U=\sqrt{p\cdot\bar\sigma}=\sqrt{E+\mathbf p\cdot\boldsymbol\sigma}$、$W=\sqrt{p\cdot\sigma}=\sqrt{E-\mathbf p\cdot\boldsymbol\sigma}$.$m>0$时两矩阵正定,取唯一正Hermitian平方根;它们可同时对角化且$UW=WU=m1$.Pauli恒等式给$\sigma^2U^*\sigma^2=W$.因$\zeta_s$为实,$v_s^*=U^*\sigma^2\zeta_s$,从而
\[
 (p\cdot\sigma)u_s-m\sigma^2v_s^*
 =W^2U\zeta_s-mW\zeta_s=0,
\]
\[
 (p\cdot\sigma)v_s+m\sigma^2u_s^*
 =-mW\sigma^2\zeta_s+m\sigma^2U^*\zeta_s=0.
\]
$U$可逆,$U\zeta_\pm$线性无关,且任意$u$经第一式唯一确定$v$,所以这两种极化给完整模式基.一般动量下$u_+^\dagger u_-$未必为零;完整性依靠可逆性,不能沿用官方解中混用$\mathbf p$和$-\mathbf p$的“正交”论证.

(d) 只引入一套费米算符:
\[
 \chi(x)=\sum_s\int\frac{\dd^3p}{(2\pi)^3\sqrt{2E_p}}
 [a_s(\mathbf p)u_s(p)\e^{-\ii px}+a_s^\dagger(\mathbf p)v_s(p)\e^{\ii px}],
 \quad\{a_r(\mathbf p),a_s^\dagger(\mathbf q)\}=\delta_{rs}(2\pi)^3\delta^{(3)}(\mathbf p-\mathbf q).
\]
验证归一化可用$\sum_su_su_s^\dagger=U^2$、$\sum_sv_sv_s^\dagger=U^2$.等时把第二项的$\mathbf p$反号,$U(\mathbf p)^2+U(-\mathbf p)^2=2E_p1$,于是$\{\chi,\chi^\dagger\}$正好给空间delta;对两个未共轭场,注意二维恒等式$X\sigma^2X^T=(\det X)\sigma^2$及$\det U(p)=m$,故
\[
 \sum_su_s(p)v_s(p)^T=U(p)\sigma^2U(p)^T=m\sigma^2,\qquad
 \sum_sv_s(p)u_s(p)^T=-m\sigma^2.
\]
等时反对易子中的两个交叉项分别乘$\e^{+\ii\mathbf p\cdot(\mathbf x-\mathbf y)}$与$\e^{-\ii\mathbf p\cdot(\mathbf x-\mathbf y)}$;将第二项$\mathbf p\mapsto-\mathbf p$后逐项抵消,因此$\{\chi,\chi\}=0$.这里矩阵恒等式使用同一个$p$,不能假称$U(p)\sigma^2U(-p)^T$一般反对称.

(e) $a_s(\mathbf p)\ket0=0$,$\braket00=1$,令$\ket{\mathbf p,s}=\sqrt{2E_p}a_s^\dagger(\mathbf p)\ket0$,便有
$\braket{\mathbf p,r}{\mathbf q,s}=2E_p(2\pi)^3\delta_{rs}\delta^{(3)}(\mathbf p-\mathbf q)$.有两个自旋态,却无独立$c_s^\dagger$反粒子态;反粒子就是粒子自身.Dirac理论每动量有两粒子、两反粒子模式,Majorana只有两模式.
\end{solution}

\originalsection{原题3}
\begin{exercise}
对$N$对独立Grassmann变量及复矩阵$A$证明
\[
 \int\prod_{i=1}^N(\dd\theta_i^*\dd\theta_i)\e^{-\theta_i^*A_{ij}\theta_j}=\det A,
\]
\[
 \int\prod_{i=1}^N(\dd\theta_i^*\dd\theta_i)\theta_k\theta_l^*\e^{-\theta_i^*A_{ij}\theta_j}
 =\det A\,(A^{-1})_{kl}.
\]
第二式假设$A$可逆,并固定Berezin积分次序.
\end{exercise}
\begin{solution}
采用$\int\dd\theta\,1=0$、$\int\dd\theta\,\theta=1$,每对次序使$\int\dd\theta^*\dd\theta\,\theta\theta^*=1$;不同变量对是偶数整体,故交换整对不改符号.

指数是有限多项式.幂次大于$N$的项必有变量重复而为零,小于$N$的项无法含全$2N$个变量,积分为零;只需$N$次项.写$-\theta_i^*A_{ij}\theta_j=A_{ij}\theta_j\theta_i^*$,则
\begin{align*}
 I(A)&=\frac1{N!}\sum_{i_1,\ldots,i_N\atop j_1,\ldots,j_N}
 A_{i_1j_1}\cdots A_{i_Nj_N}
 \int\prod_r(\dd\theta_r^*\dd\theta_r)
 (\theta_{j_1}\theta_{i_1}^*)\cdots(\theta_{j_N}\theta_{i_N}^*)\\
 &=\frac1{N!}\epsilon_{i_1\cdots i_N}\epsilon_{j_1\cdots j_N}
 A_{i_1j_1}\cdots A_{i_Nj_N}=\det A.
\end{align*}
两个epsilon来自分别排列两种变量;把所有未星变量移到左边引入$(-1)^{N(N-1)/2}$,重组为标准变量对又引入同一个符号,互相抵消.最后一式是行列式的Leibniz公式.

对矩阵元素$A_{lk}$求普通导数,有限多项式允许逐项求导,$\partial I/\partial A_{lk}$的被积函数多出$-\theta_l^*\theta_k=\theta_k\theta_l^*$.行列式求导由$\det(A+\delta A)=\det A\det(1+A^{-1}\delta A)$及$\det(1+\epsilon B)=1+\epsilon\tr B+\order(\epsilon^2)$给出
$\partial\det A/\partial A_{lk}=\det A(A^{-1})_{kl}$,即第二式.奇异$A$时逆矩阵无定义,但右端可用相应余子式延拓为多项式.
\end{solution}

\originalsection{原题4}
\begin{exercise}
考虑原作业的Yukawa理论
$\Lag_{\rm o}=-\frac12(\widehat\partial\phi)^2-\frac12m^2\phi^2-\ii\bar\psi_{\rm o}(\Gamma\cdot\partial-m)\psi-g\phi\bar\psi_{\rm o}\psi$.
标量用虚线,费米子用有向实线,称粒子为$p$、反粒子为$\bar p$.画各过程最低阶图并写振幅:
(a) $\bar p+\bar p\to\bar p+\bar p$,初态$(p_1,s_1),(p_2,s_2)$,末态$(p'_1,s'_1),(p'_2,s'_2)$;
(b) $p+\bar p\to\phi+\phi$,初态同上,末态$p'_1,p'_2$;
(c) $p+\phi\to p+\phi$,初态$(p_1,s_1),p_2$,末态$(p'_1,s'_1),p'_2$.
\end{exercise}
\begin{solution}
先固定振幅约定:$S_{fi}=\delta_{fi}+\ii(2\pi)^4\delta^{(4)}(P_f-P_i)\mathcal M_{fi}$.官方解把图的完整$\ii\mathcal M$也记作$M$,本解明确区分.
原$-g\phi\bar\psi_{\rm o}\psi=-\ii g\phi\bar\psi\psi$,故统一写$-y\phi\bar\psi\psi$,$y=\ii g$.若要求Hermitian物理作用量,$y$应实,即原$g=-\ii y$;原题的耦合符号没有展开这项相位约定.下列计算按本书实$y$进行,代回$y=\ii g$恢复原形式.为区分内线质量,记$m_\psi,m_\phi$,原题二者均为$m$.

顶点$-\ii y$,标量传播子$\ii/(q^2-m_\phi^2+\ii0)$,费米传播子$\ii(\slashed q+m_\psi)/(q^2-m_\psi^2+\ii0)$.下面三行依次是(a)、(b)、(c);实线箭头只表费米流,虚线为标量,(a)右图交换末态腿,(b)右图交换两个末态标量,(c)右图为交叉道.
\begin{center}\begin{tikzpicture}[>=Stealth,scale=.8]
\foreach \x/\lab in {0/t,5/u}{
\draw[->](\x+1,1)--(\x-1,1);\draw[->](\x+1,-1)--(\x-1,-1);\draw[dashed](\x,1)--(\x,-1);\node at(\x,-1.6){(a) $\lab$};}
\foreach \x/\lab in {0/t,5/u}{
\draw[->](\x-1,-3)--(\x,-3.5);\draw[->](\x,-3.5)--(\x,-4.5);\draw[->](\x,-4.5)--(\x-1,-5);
\draw[dashed](\x,-3.5)--(\x+1,-3);\draw[dashed](\x,-4.5)--(\x+1,-5);\node at(\x,-5.6){(b) $\lab$};}
\draw[->](-1,-8)--(0,-7.5);\draw[->](0,-7.5)--(1,-7.5);\draw[->](1,-7.5)--(2,-8);
\draw[dashed](-1,-7)--(0,-7.5);\draw[dashed](1,-7.5)--(2,-7);\node at(.5,-8.6){(c) $s$};
\draw[->](4,-8)--(5,-7.5);\draw[->](5,-7.5)--(6,-7.5);\draw[->](6,-7.5)--(7,-8);
\draw[dashed](4,-7)--(6,-7.5);\draw[dashed](5,-7.5)--(7,-7);\node at(5.5,-8.6){(c) $u$};
\end{tikzpicture}\end{center}
(a) 有$t=(p_1-p'_1)^2$、$u=(p_1-p'_2)^2$两图.固定外部态次序后,交换两个同种费米子多一个负号:
\[
 \mathcal M_a=-y^2\left[
 \frac{(\bar v_{s_1}(p_1)v_{s'_1}(p'_1))(\bar v_{s_2}(p_2)v_{s'_2}(p'_2))}{t-m_\phi^2+\ii0}
 -\frac{(\bar v_{s_1}(p_1)v_{s'_2}(p'_2))(\bar v_{s_2}(p_2)v_{s'_1}(p'_1))}{u-m_\phi^2+\ii0}\right].
\]
原图若用交叉外腿画$u$道,应同时交换两个末态标签;只有拓扑相同不表示两个图相同.

(b) 两个标量相同,但交换它们不产生费米态交换负号.沿唯一费米线的动量分别是$q_t=p_1-p'_1$、$q_u=p_1-p'_2$:
\[
 \mathcal M_b=-y^2\bar v_{s_2}(p_2)
 \left[\frac{\slashed q_t+m_\psi}{t-m_\psi^2+\ii0}
 +\frac{\slashed q_u+m_\psi}{u-m_\psi^2+\ii0}\right]u_{s_1}(p_1).
\]
$c$型反粒子外线收缩给$\bar v$,粒子入射外线给$u$;矩阵顺序必须沿线保留.

(c) 费米子可先吸收入射标量,也可先发射出射标量,故有$s=(p_1+p_2)^2$和$u=(p_1-p'_2)^2$两图:
\[
 \mathcal M_c=-y^2\bar u_{s'_1}(p'_1)
 \left[\frac{\slashed p_1+\slashed p_2+m_\psi}{s-m_\psi^2+\ii0}
 +\frac{\slashed p_1-\slashed p'_2+m_\psi}{u-m_\psi^2+\ii0}\right]u_{s_1}(p_1).
\]
每式由两个$-\ii y$顶点乘一个$\ii$传播子再除定义中的$\ii$得到;此处只有开费米线,没有闭环负号.
\end{solution}
